\documentclass[runningheads]{llncs}
\usepackage[T1]{fontenc}
\usepackage{graphicx}
\usepackage{amsmath}
\usepackage{booktabs}

% ============================================================================
% ACIIDS 2027 -- 19th Asian Conference on Intelligent Information and
% Database Systems. Bali, Indonezja, 5-7 kwietnia 2027.
%
% SPRAWDZONE NA STRONIE KONFERENCJI (aciids.pwr.edu.pl/2027, 8.10.2026):
%   deadline       31.10.2026, oznaczony "Extended" (pierwotnie 31.07, potem 30.09)
%   seria          LNCS / LNAI (Springer)  -> klasa llncs, ta sama co tutaj
%   dlugosc        12-15 stron WLACZNIE z rysunkami, tabelami i bibliografia
%                  (powyzej doplata, absolutne maksimum 25 stron)
%   zgloszenie     PDF przez EasyChair: easychair.org/conferences/?conf=aciids2027
%   recenzja       dwustopniowa; strona NIE mowi, czy slepa -> zostawiam nazwiska,
%                  do potwierdzenia przed wyslaniem
%   uwaga          prace sa przepuszczane przez kontrole podobienstwa, a
%                  autoplagiat = desk reject. Mamy dwie prace o pokrywajacej sie
%                  tematyce (windowed_qubo, qaoa_gcmm), wiec sekcje 2-3 musza byc
%                  pisane od nowa, nie przeklejane. W GCMM zmierzylismy 0,61%
%                  pokrycia -- utrzymac ten poziom.
%
% SZKIELET -- sekcje, po 2-3 zdania, tabele tam gdzie sa dane.
% Status danych: [MAM] gotowe, [TRANSZA] po transzy 2000 kredytow, [BRAK] do zdobycia.
% ============================================================================

\begin{document}

\title{Windowed QAOA Neighborhoods for the Permutation Flow Shop Problem
       on the Full Taillard Benchmark}
\author{Remigiusz Wojew\'{o}dzki \and Wojciech Bo\.{z}ejko}
\authorrunning{R. Wojew\'{o}dzki and W. Bo\.{z}ejko}
\institute{Wroc\l{}aw University of Science and Technology, Faculty of
  Information and Communication Technology, Department of Control and
  Quantum Computing, Wroc\l{}aw, Poland\\
  \email{wojciech.bozejko@pwr.edu.pl}}
\maketitle

\begin{abstract}
% [PISAC NA KONCU] 150-250 slow. Teza: dekompozycja okienkowa przenosi cztery
% sasiedztwa bramkowe na caly zakres Taillarda na 16 kubitach; jedna ewaluacja
% sasiedztwa przekracza najdluzszy budzet protokolu, wiec porownanie idzie na
% ruch, nie na sekunde; PRD lezy wysoko nad klasyka i annealingiem.
\keywords{Permutation flow shop \and QAOA \and Quantum neighborhoods
          \and Iterated local search \and Taillard benchmark}
\end{abstract}

% =====================================================================
\section{Introduction}
% [MAM] -- proza ciagla, bez wypunktowan
The permutation flow shop problem remains strongly NP-hard for three or more
machines, and practical solvers evaluate neighborhoods millions of times inside
local search. Recent work replaces that evaluation with a quantum subroutine,
either on an annealer or on a gate processor, but every gate-model study so far
stops at twenty jobs because the swap-selection QUBO needs $K = n-1$ variables
for the linear neighborhoods and $O(n^2)$ for the dense ones, while current
devices hold sixteen to fifty-four qubits.

This paper removes that ceiling. A windowed decomposition carries all four
neighborhoods to the entire Taillard range on a sixteen-qubit processor, and the
measurements that follow are the first for gate-model neighborhoods beyond
$n = 20$. The results also show where the comparison protocol inherited from the
annealing line breaks down, since one neighborhood evaluation already costs more
wall-clock time than the longest budget under which the classical baselines were
measured.

% =====================================================================
\section{Related Work}
% [MAM] -- skrotowo, 1 akapit
% Taillard 1993; Smutnicki blocks; Wodecki NPI; QUBO dla PFSP (Bozejko, Venturelli,
% Amaro); QAOA fixed-angle (Farhi, Wurtz); nasze dwa poprzednie artykuly.

% =====================================================================
\section{Problem and QUBO Formulation}
\label{sec:method}

\subsection{Permutation flow shop and the four neighborhoods}
\label{subsec:pfsp}
An instance gives $n$ jobs and $m$ machines, each job passing the machines in
the same order, and a solution is a permutation $\pi$ of the jobs whose
makespan $C_{\max}(\pi)$ is the moment the final job clears the final
machine. The four neighborhoods differ in how far a single move reaches along
$\pi$: Adjacent exchanges one pair of consecutive jobs, Fibonacci applies a set
of pairwise disjoint consecutive exchanges, and Dynasearch and Motzkin select
non-crossing intervals whose endpoints are exchanged. That reach is what later
fixes how wide a window may be, so we record it here rather than in the
decomposition.

For every candidate move the makespan change $\delta_k$ follows in constant
time from the Head and Tail profiles of the current permutation, so a whole
candidate set is priced in one pass. Two classical filters cut that set before
it ever reaches a solver: the block property of Smutnicki discards moves that
cannot touch the critical path, and the no-positive-improvement rule drops
moves with $\delta_k \ge 0$. Both matter for what follows, because they make
the number of surviving candidates depend on the instance rather than on $n$
alone, and that number is what a device has to hold.

\subsection{Swap-selection QUBO}
\label{subsec:qubo}
Let the neighborhood supply $K$ candidates and let $x_k \in \{0,1\}$ record
whether candidate $k$ is applied, for $k = 1, \dots, K$. Choosing a compatible
subset is the quadratic program
\begin{equation}
  \min_{\mathbf{x} \in \{0,1\}^{K}} \mathbf{x}^{\mathsf T} Q \mathbf{x},
  \qquad Q_{kk} = \delta_k, \qquad
  Q_{kl} = \lambda \ \text{ for } 1 \le k < l \le K \text{ in conflict},
  \label{eq:qubo}
\end{equation}
with $Q$ upper triangular so that each conflicting pair is charged once, and
with $\lambda = \sum_{k=1}^{K} \lvert \delta_k \rvert + 1$ large enough that no
infeasible selection can beat a feasible one. Adjacent is the one exception:
there every pair conflicts, so the constraint becomes one-hot and the penalty
moves onto the diagonal as well, with $Q_{kk} = \delta_k - P$ and
$Q_{kl} = 2P$, where $P = 2\max_k \lvert \delta_k \rvert + 1$.

Which pairs are charged is the only thing that separates the four
neighborhoods, and it is also what decides whether the matrix survives being
cut down to a window. Fibonacci charges a pair only when the two exchanges
touch the same position, which makes $Q$ tridiagonal, and Adjacent charges
every pair. Both patterns depend on the candidates alone, so restricting them
to a subset of positions yields exactly the matrix the full problem would have
assigned to that subset. The interval neighborhoods charge a pair whenever two
intervals cross, which is a relation between candidates rather than positions,
and Section~\ref{sec:windowed} treats them separately for that reason.

\subsection{Fixed-angle QAOA}
\label{subsec:qaoa}
Substituting $x_k = (Z_k + 1)/2$ turns~\eqref{eq:qubo} into an Ising
Hamiltonian $H_C$ over $K$ spins, and QAOA prepares
$\lvert \psi \rangle = e^{-i\beta H_M} e^{-i\gamma H_C} H^{\otimes K}
\lvert 0 \rangle^{K}$ at depth $p = 1$, with the transverse field
$H_M = \sum_{k=1}^{K} X_k$ as mixer and $\gamma \in [0, \pi]$,
$\beta \in [0, \pi/2]$. The two angles are optimized once per neighborhood on
a small set of representative matrices and then frozen, which keeps the online
cost of a move at one submission instead of the hundreds of circuits a
variational loop would spend before committing to a single move.

A move samples the circuit and keeps whichever assignment scores lowest
against the original $Q$, rather than whichever one came up most often. The
distinction decides whether the method works at all on a noisy device: at
$p = 1$ the modal outcome is usually the all-zeros assignment, which applies no
swap, or a high-energy string with many bits set. Energy-based readout also
makes the windowed solver of the next section well defined, since each window
is sampled and reduced independently and the winner is whichever assignment its
own matrix scores lowest.

% =====================================================================
\section{Windowed Decomposition}
\label{sec:windowed}

\subsection{Why the full QUBO does not fit}
% [MAM] -- Tabela 1
The width of the linear QUBO is instance-dependent, because the block
accelerator removes candidates outside the critical block. Table~\ref{tab:cap}
reports it across all 120 Taillard instances against the capacity of the three
IQM processors, and shows that without decomposition the linear neighborhoods
reach 24 of 120 instances on sixteen qubits and 79 on fifty-four.

\begin{table}[t]
\caption{Width of the undecomposed linear QUBO over the ten instances of each
configuration, and the number of instances that fit each device.}
\label{tab:cap}
\centering
\begin{tabular}{@{}lrrrrrrrr@{}}
\toprule
& & & \multicolumn{3}{c}{variables} & \multicolumn{3}{c}{instances fitting} \\
\cmidrule(lr){4-6}\cmidrule(lr){7-9}
file & $n$ & $m$ & min & med & max & 16q & 20q & 54q \\
\midrule
tai20\_5 & 20 & 5 & 5 & 10 & 18 & 9 & 10 & 10 \\
tai20\_10 & 20 & 10 & 11 & 15 & 19 & 8 & 10 & 10 \\
tai20\_20 & 20 & 20 & 11 & 16 & 18 & 5 & 10 & 10 \\
tai50\_5 & 50 & 5 & 10 & 24 & 32 & 2 & 3 & 10 \\
tai50\_10 & 50 & 10 & 19 & 30 & 38 & 0 & 1 & 10 \\
tai50\_20 & 50 & 20 & 26 & 36 & 42 & 0 & 0 & 10 \\
tai100\_5 & 100 & 5 & 17 & 30 & 52 & 0 & 1 & 10 \\
tai100\_10 & 100 & 10 & 31 & 48 & 71 & 0 & 0 & 7 \\
tai100\_20 & 100 & 20 & 49 & 65 & 78 & 0 & 0 & 1 \\
tai200\_10 & 200 & 10 & 42 & 70 & 93 & 0 & 0 & 1 \\
tai200\_20 & 200 & 20 & 77 & 101 & 117 & 0 & 0 & 0 \\
tai500\_20 & 500 & 20 & 129 & 184 & 219 & 0 & 0 & 0 \\
\midrule
total & & & & & & 24 & 35 & 79 \\
\bottomrule
\end{tabular}
\end{table}

\subsection{Window geometry}
% [MAM] -- kluczowe: dwa rezimy skalowania
A window of $w$ consecutive positions yields $w-1$ variables for the
tridiagonal and one-hot matrices but $w(w-1)/2$ for the pairwise one, so one
window size cannot serve both classes. On sixteen qubits this means $w = 17$ for
the linear neighborhoods against $w = 6$ for the dense ones, and the window of
19 used by the annealing study fits no gate device, since it would need 171
variables.

\subsection{Tridiagonal and pairwise windows}
% [MAM] -- fibonacci nowe, dynasearch/motzkin istnialo
Fibonacci's QUBO couples only physically adjacent swaps, so the matrix of a
position window is exactly the global matrix restricted to that window and the
decomposition leaves the formulation unchanged. Swaps straddling a window
boundary are resolved by the same non-overlap validation the single-QUBO path
uses.

\subsection{One-hot windows}
% [MAM] -- adjacent, dwupoziomowy argmin; uzasadnienie ze to ten sam cel
The global one-hot QUBO selects the candidate of lowest delta, so solving one
one-hot QUBO per window and keeping the best proposal computes the same
$\arg\min$ and still applies exactly one swap. This matters because the D-Wave
enhanced variants never decomposed the linear neighborhoods; they relied on the
annealer holding around 190 logical variables, which no gate device does.

\subsection{Batching the windows of a move}
% [MAM]
Every window of one move is submitted as a single job, so the queue and the job
overhead are paid once per move rather than once per window.

% =====================================================================
\section{Experimental Protocol}
\label{sec:protocol}

\subsection{Benchmarks, metaheuristic and metrics}
% [MAM] -- 12 konfiguracji, 5 instancji, ILS, 10 s, p=1, 4096 shotow, PRD od LB
We follow the protocol of the windowed annealing study so that the rows stay
comparable: Taillard instances with $n \in \{20, 50, 100, 200, 500\}$ and
$m \in \{5, 10, 20\}$, percentage relative deviation from the Taillard lower
bound, and a wall-clock budget of ten seconds. Each neighborhood drives an
iterated-local-search loop on the \texttt{sirius} processor, sixteen qubits in a
star topology, with $p = 1$ and 4096 shots per circuit.

\subsection{Deviations from the annealing protocol}
% [MAM] -- trzy odstepstwa, kazde z uzasadnieniem
Three departures are deliberate. Only iterated local search is run, because a
single move from the identity permutation cannot distinguish it from simulated
annealing; one run replaces five seeds, because repeated runs sample shot noise
alone and the measured spread is 0.00 points for Adjacent and 0.03 for
Fibonacci; and five instances replace ten, because at the measured rate of
0.75 credits per circuit the full ten-instance protocol would cost about 12\,750
credits.

% =====================================================================
\section{Results}
\label{sec:results}

\subsection{Solution quality across the Taillard range}
% [TRANSZA] -- Tabela 2 jest rdzeniem artykulu
% Po transzy: 12 konfiguracji x 4 sasiedztwa, srednie PRD po 5 instancjach.
% Ponizej stan obecny: instancja 0, piec plikow -- do zastapienia.

\begin{table}[t]
\caption{Mean PRD [\%] on \texttt{sirius}. PRELIMINARY: instance 0 only, the
dense columns averaged over five repetitions. To be replaced by five instances
per configuration.}
\label{tab:rpd}
\centering
\begin{tabular}{@{}lrrrrrrr@{}}
\toprule
file & $n$ & $m$ & Adjacent & Fibonacci & Dynasearch & Motzkin & moves \\
\midrule
tai20\_5 & 20 & 5 & 12.34 & 9.17 & 13.08 & 15.58 & 2 \\
tai20\_10 & 20 & 10 & 34.12 & 30.18 & 23.95 & 26.08 & 2 \\
tai20\_20 & 20 & 20 & 40.19 & 37.89 & 41.39 & 44.95 & 2 \\
tai50\_5 & 50 & 5 & 11.95 & 12.24 & 11.59 & 11.59 & 1 \\
tai50\_10 & 50 & 10 & 28.07 & 26.87 & 24.10 & -- & 1 \\
tai50\_20 & 50 & 20 & -- & -- & -- & -- & -- \\
tai100\_5 & 100 & 5 & -- & -- & -- & -- & -- \\
tai100\_10 & 100 & 10 & -- & -- & -- & -- & -- \\
tai100\_20 & 100 & 20 & -- & -- & -- & -- & -- \\
tai200\_10 & 200 & 10 & -- & -- & -- & -- & -- \\
tai200\_20 & 200 & 20 & -- & -- & -- & -- & -- \\
tai500\_20 & 500 & 20 & -- & -- & -- & -- & -- \\
\bottomrule
\end{tabular}
\end{table}

\subsection{Comparison with the classical and annealing baselines}
% [TRANSZA dla kolumny bramkowej; klasyczne i annealingowe MAM w results/]
% Tabela 3: te same instancje, trzy klasy. Klasyczny ILS na tai20_5 przy 10 s
% daje 3.41% dla Adjacent wobec 12.34% bramkowego -- rzad wielkosci.

\subsection{Collapse of the time-budget axis}
% [MAM czesciowo] -- Tabela 4
One move takes between three and seventy-three seconds depending on $n$ and on
the neighborhood class, so the ten-second budget admits two moves at $n = 20$
and one from $n = 50$ upward. Every shorter budget of the protocol therefore
describes the same computation, which is why the comparison is reported per move
rather than per second.

\subsection{Resource model}
% [MAM] -- Tabela 5, w calosci zmierzona
Billing follows QPU time and therefore circuit count: the job log gives 1.5
seconds per circuit, which at the published rate is about 0.75 credits per
circuit. Table~\ref{tab:res} turns that into the cost of one move for each
configuration, and shows that the two largest size classes account for roughly
half of the full protocol.

\begin{table}[t]
\caption{Circuits per move, QPU seconds and credits per move, with windows of 17
for the linear neighborhoods and 6 for the dense ones. Measured on
\texttt{sirius}, averaged over the ten instances of each configuration.}
\label{tab:res}
\centering
\begin{tabular}{@{}lrrrrrrrr@{}}
\toprule
& & & \multicolumn{2}{c}{circuits} & \multicolumn{2}{c}{seconds}
& \multicolumn{2}{c}{credits} \\
\cmidrule(lr){4-5}\cmidrule(lr){6-7}\cmidrule(lr){8-9}
file & $n$ & $m$ & lin. & dense & lin. & dense & lin. & dense \\
\midrule
tai20\_5 & 20 & 5 & 2.0 & 3.5 & 3 & 5 & 1.5 & 2.6 \\
tai20\_10 & 20 & 10 & 2.0 & 4.7 & 3 & 7 & 1.5 & 3.5 \\
tai20\_20 & 20 & 20 & 2.0 & 5.7 & 3 & 9 & 1.5 & 4.3 \\
tai50\_5 & 50 & 5 & 5.6 & 4.8 & 8 & 7 & 4.2 & 3.6 \\
tai50\_10 & 50 & 10 & 5.9 & 8.4 & 9 & 13 & 4.4 & 6.3 \\
tai50\_20 & 50 & 20 & 6.0 & 11.1 & 9 & 17 & 4.5 & 8.3 \\
tai100\_5 & 100 & 5 & 8.6 & 6.2 & 13 & 9 & 6.4 & 4.7 \\
tai100\_10 & 100 & 10 & 10.5 & 10.4 & 16 & 16 & 7.9 & 7.8 \\
tai100\_20 & 100 & 20 & 11.8 & 18.8 & 18 & 28 & 8.9 & 14.1 \\
tai200\_10 & 200 & 10 & 18.5 & 14.1 & 28 & 21 & 13.9 & 10.6 \\
tai200\_20 & 200 & 20 & 21.6 & 25.2 & 32 & 38 & 16.2 & 18.9 \\
tai500\_20 & 500 & 20 & 48.7 & 28.4 & 73 & 43 & 36.5 & 21.3 \\
\bottomrule
\end{tabular}
\end{table}

% =====================================================================
\section{Discussion}
\label{sec:discussion}

\subsection{The first move agrees across coupling maps}
% [MAM] -- mocny wynik, potwierdza teze progowa z GCMM na trzeciej architekturze
Adjacent returns a makespan of 1390 on instance 0 of \texttt{tai20\_5} on
\texttt{garnet}, \texttt{emerald} and \texttt{sirius} alike, three different
coupling maps. The processors differ only in how many moves fit the budget, so a
cross-architecture comparison that does not hold the move count fixed measures
queue latency rather than fidelity.

\subsection{Non-additive deltas in the dense neighborhoods}
% [MAM] -- 8 z 52 przebiegow bez poprawy, motzkin na tai20_20 piec razy z pieciu
Eight of the fifty-two hardware runs returned the starting makespan, and Motzkin
on \texttt{tai20\_20} did so in all five repetitions. The deltas entering the
QUBO are computed independently, while the pairwise penalties do not capture how
several swaps compose, so a set of individually improving swaps can worsen the
schedule; this limits the formulation rather than the device, and it applies to
the annealing line as well.

\subsection{What a full campaign would cost}
% [MAM]
% Pelny protokol windowed_qubo (10 instancji, 5 powtorzen dla gestych) ~12 750
% kredytow; z jednym powtorzeniem ~4 200; piec instancji ~2 100.

% =====================================================================
\section{Reproducibility}
\label{sec:repro}
% [MAM] -- tryb awarii i straz; w literaturze QPU nikt tego nie opisuje
A declined job returns an empty list of histograms without raising, and each
neighborhood then falls back to a classical greedy swap, so the run completes
and is stored as a hardware measurement while being entirely classical. We
detected this on a motzkin row whose moves took 2.5 seconds against the 12
seconds of the valid rows, and the guard now compares the number of returned
histograms against the number of submitted circuits.

% =====================================================================
\section{Conclusions}
\label{sec:conclusions}
% [NA KONIEC] -- max 2 akapity (uwaga promotora z GCMM)
% Akapit 1: metoda dosięga n=500 na 16 kubitach, cztery sasiedztwa, pierwsze
% pomiary poza n=20.
% Akapit 2: bariera jest zasobowa i policzona; czego wymagalaby pelna kampania.

\begin{thebibliography}{8}
% [DO UZUPELNIENIA] -- przeniesc z qaoa_gcmm te, ktore sie powtarzaja
\bibitem{taillard1993}
Taillard, E.: Benchmarks for basic scheduling problems.
Eur. J. Oper. Res. \textbf{64}(2), 278--285 (1993)
\end{thebibliography}

\end{document}
